\section{Preliminaries}

\subsection{点云表征}
点云以一组三维点表示工件表面，每个点除坐标外还可携带温度、应力等物理量，记为 $X=\{x_i\}_{i=1}^{N}$，$x_i\in\mathbb{R}^{C}$。点云是无序集合，同一形状对应 $N!$ 种排列，故处理点云的模型需具备置换不变性~\cite{qi2017pointnet}。与规则网格相比，点云不受拓扑与分辨率限制，可直接表示任意几何形状，适合承载有限元仿真输出的几何与场量。

\subsection{连续时间隐状态建模}
Neural ODE 以参数化的向量场描述状态的连续演化：给定初值 $c(s_0)=c_0$，任意位置的状态由
\begin{equation}
\frac{\mathrm{d}c}{\mathrm{d}s}=g(c,s;\theta)
\end{equation}
积分得到~\cite{chen2018neural}。与离散时间递推相比，它以连续自变量刻画动力学，可在任意时刻求值，并可用伴随法反向传播。这类模型已在工业过程的连续时间建模中得到应用~\cite{yuan2022autonomous, yuan2023ode, yuan2022differential}。

\subsection{世界模型}
世界模型将高维观测压缩为低维隐状态，并在隐空间中学习状态随动作的演化，从而在不与真实环境交互的条件下预测未来。典型结构含编码器、推理器与解码器三部分：编码器由观测得到隐状态，推理器预测下一时刻的隐状态，解码器将隐状态还原为观测。预测在低维空间完成，计算代价远低于直接在高维观测上推演。

\subsection{DeepONet}
DeepONet 用于学习函数到函数的映射，由 branch 与 trunk 两个子网构成：branch 接收输入函数的离散采样，trunk 接收输出函数的自变量，二者输出做内积即得该自变量处的函数值~\cite{lu2021deeponet}。由于输出是自变量的连续函数，模型可在任意位置给出预测，无需预先固定离散网格。
